% Paper skeleton for the UMR-only rewrite (see notes/outline-umr.md for the advisor's feedback).
% Every section and paragraph holds a summary of what it will say, with the numbers it will use; not final prose.
% Status tags: [data] numbers exist; [analysis] raw records exist, analysis to do;
% [new run] experiment to run; [verify] re-check by hand or with the official scorers.
% New numbers (Smatch, layers, edit coding) come from SkillParseUMR experiment/annoforge/af/{smatch_score,layers}.py,
% copied to data/annoforge/{smatch.json,layers.json,edit_codes.tsv}.
% Compile on its own: latexmk -pdf outline_umr.tex (uses acl.sty and references.bib of this folder).
\documentclass[11pt]{article}
\usepackage[review]{acl}
\usepackage{times}
\usepackage{latexsym}
\usepackage[T1]{fontenc}
\usepackage[utf8]{inputenc}
\usepackage{microtype}
\usepackage{inconsolata}
\usepackage{graphicx}
\usepackage{amsmath}
\usepackage{booktabs}
\usepackage{multirow}
\usepackage{array}
\usepackage{enumitem}
\usepackage{xspace}
\usepackage{xcolor}
\usepackage{CJKutf8}
\usepackage{tipa}

\newcommand{\jumach}{ju\textlengthmark{}m\ae{}\textteshlig{}\xspace}
\newcommand{\umr}{\textsc{UMR}\xspace}
\newcommand{\zh}[1]{\begin{CJK}{UTF8}{gbsn}#1\end{CJK}}
\newcommand{\stag}[2]{\textcolor{#1}{\textsf{\scriptsize[#2]}}}
\newcommand{\HAVE}{\stag{green!50!black}{data}\xspace}
\newcommand{\ANA}{\stag{blue!70!black}{analysis}\xspace}
\newcommand{\RUN}{\stag{red!75!black}{new run}\xspace}
\newcommand{\CHK}{\stag{orange!80!black}{verify}\xspace}
\newcommand{\tbd}{\textcolor{red}{TBD}\xspace}
\newcommand{\placeholder}[2]{\fbox{\parbox{#1}{\centering\small\vspace{0.6em}#2\vspace{0.6em}}}}
\newcommand{\ex}[1]{{\ttfamily\scriptsize #1}}

\title{Learning to Parse Uniform Meaning Representation\\
from Its Guideline and Fifty Sentences}

\author{Anonymous submission}

\begin{document}
\maketitle

% ==================================================================
\begin{abstract}
\emph{Five sentences.}
(1) Uniform Meaning Representation (UMR) unifies concept graphs, PropBank senses, aspect, reference attributes and document-level temporal, modal and coreference relations across languages, but most languages have little annotated data, and the fine-tuned systems of the first UMR parsing shared task reach 37.7 (Chinese) and 40.3 (English) AnCast++ F on our test documents.
(2) An LLM given the 366K-character guideline does no better: it scores zero on document-level modality and near zero on temporal relations and coreference, because the guideline writes UMR differently from the data.
(3) We let an LLM revise a step-wise annotator on 50 gold sentences, editing its instructions, writing a PropBank lookup and adding rule code, which raises Chinese and English F by 6.1 and 11.3, Smatch by 6.0 and 9.9, and the official \jumach score to 42.9 and 38.0, above the fine-tuned systems (38.8 and 27.3).
(4) Tracing every edit to the layer it targets, we find that edits stating a convention carry over to unseen documents almost fully, edits giving the model a lexical resource carry over in part, and edits stating a semantic decision criterion (aspect, coreference, temporal anchoring) improve the training sentences but not the test documents.
(5) The layers that do not move are not those on which annotators disagree: two annotations of the same English text agree at 83--92 on aspect, temporal relations and coreference, where the model stays at 13--60.
\end{abstract}

% ==================================================================
\section{Introduction}
\label{sec:intro}

\paragraph{UMR and why it matters}
UMR \citep{vangysel2021designing,bonn2024building} extends AMR \citep{banarescu2013abstract} with aspect, modality, person and number, word alignments, and document-level temporal, modal and coreference relations, and is designed for many languages, including those without a PropBank. One representation across sentences and languages makes it a demanding but useful parsing target. \HAVE

\paragraph{A complex scheme with little data}
The guideline runs to 366K characters and needs external resources (Chinese and English PropBank; \citealp{xue2009adding,palmer2005propbank}). Gold data are scarce: the released English data are mostly partial conversions from AMR, and only three English documents have full sentence- and document-level annotation. The fine-tuned systems of the first shared task \citep{stepanek2026first,heinecke2026orange,devergnette2026sema} reach 37.7 / 40.3 F and 54.2 / 56.5 Smatch on our documents (Chinese / English). \HAVE

\paragraph{Reading the guideline is not enough}
Haiku with the whole guideline in the prompt scores 30.9 F and 12.5 \jumach on Chinese, and 0 on document-level modality in both languages (\S\ref{sec:diagnosis}). The guideline writes UMR in a notation that the released data do not use (\texttt{:modstr FullAff}, \texttt{AUTH}, \texttt{DCT}, \texttt{:ref-number Singular}); what the model lacks is the project's conventions and its lexical resources, not an understanding of the sentence. \HAVE

\paragraph{Approach}
A builder LLM writes a three-step annotator from the guideline, then revises it over ten rounds of five gold sentences, guided by a judge that reports errors per UMR layer: it edits the step instructions, writes a PropBank lookup, and adds rule code. Unlike guideline refinement for flat labels \citep{kim2026refining,kim2026diziner}, the target is a multi-layer graph and the edits reach instructions, tools and code. \HAVE

\paragraph{Results}
+6.1 / +11.3 F and +6.0 / +9.9 Smatch over the seed; \jumach 42.9 / 38.0 against 38.8 / 27.3 for the best fine-tuned system; the full Chinese gain is reached after three rounds (15 sentences). \HAVE

\paragraph{What we learn}
Coding all 168 applied edits by type and target layer (\S\ref{sec:analysis}) shows three regimes: conventions transfer, lexical lookups transfer in part, and decision criteria for aspect, temporal anchoring and coreference fit the training group without improving the test documents. Human agreement on these layers is high, so the limit is in what fifty sentences can teach, not in noisy gold. \HAVE{} \CHK

\paragraph{Contributions}
\begin{enumerate}[leftmargin=*,itemsep=1pt,topsep=2pt]
  \item A per-layer diagnosis of LLM UMR parsing from the guideline (\S\ref{sec:diagnosis}).
  \item An annotator learned from 50 gold sentences, without fine-tuning, that outperforms the fine-tuned shared-task systems on every metric (\S\ref{sec:main}).
  \item An account of every edit and its effect, layer by layer, explaining which parts of UMR can be learned from a few examples and which cannot (\S\ref{sec:analysis}).
\end{enumerate}

% ==================================================================
\section{UMR and Related Work}
\label{sec:umr}

\subsection{What UMR represents}

\paragraph{What UMR adds to AMR}
AMR is English-centric, has no tense or aspect, little modality, single sentences only, and no word alignment. UMR adds a cross-lingual design (predicates without a frame use \texttt{lemma-00} with \texttt{:actor}/\texttt{:undergoer}), aspect, modal strength as a document-level modal dependency, temporal dependency, cross-sentence coreference, reference attributes and alignment \citep{vangysel2021designing,bonn2025uniform,bonn2026uniform}. Figure~\ref{fig:example}. \HAVE

\begin{figure}[t]
\centering
\placeholder{0.95\columnwidth}{Figure 1: a Chinese test sentence with its sentence graph, alignment and document-level \texttt{:temporal}, \texttt{:modal}, \texttt{:coref} groups, each part labelled with its layer of Table~\ref{tab:layers}.}
\caption{A UMR annotation and its layers. \ANA}
\label{fig:example}
\end{figure}

\paragraph{Eleven layers and four kinds of difficulty}
Table~\ref{tab:layers} groups the 25 checks of our judge into eleven layers and assigns each a kind of difficulty: \emph{notation} (how the project writes a decision), \emph{resource} (a lexicon must be consulted), \emph{discourse} (a judgment over the sentence and its context), \emph{format} (well-formed document-level groups, variables and alignments). This table organizes the analysis. \HAVE

\begin{table}[t]
\centering\small
\setlength{\tabcolsep}{3pt}
\resizebox{\columnwidth}{!}{%
\begin{tabular}{@{}lllll@{}}
\toprule
Layer & & Example decision & Difficulty & Checks \\
\midrule
Concepts & CON & event or entity; Chinese word kept & notation & 5 \\
Senses & SEN & \texttt{lemma-01} vs.\ \texttt{-02}; \texttt{-00} & resource & 1 \\
Arguments & ARG & ARG0 vs.\ ARG1; \texttt{:place} vs.\ \texttt{:location} & resource / notation & 6 \\
Aspect & ASP & state, habitual, performance & discourse & 1 \\
Other attributes & ATT & polarity, refer-person/number & notation & 3 \\
Names & NAM & name tokens, wiki & notation & 2 \\
Reentrancy & REE & pronoun reuses a variable & notation / discourse & 1 \\
Doc.\ temporal & TMP & which events anchor to DCT & discourse & 1 \\
Doc.\ modal & MOD & \texttt{(author :full-affirmative e)} & notation & 1 \\
Coreference & COR & same-entity, subset-of & discourse & 2 \\
Doc.\ format & FMT & groups, variables, alignment & format & 1 \\
\bottomrule
\end{tabular}}
\caption{UMR layers, the kind of difficulty each poses, and the number of judge checks pooled into each.}
\label{tab:layers}
\end{table}

\subsection{Data and prior systems}

\paragraph{Shared task}
The first UMR parsing shared task \citep{stepanek2026first} on UMR~2.2 \citep{bonn2026uniform}; its systems Orange \citep{heinecke2026orange} and Sema \citep{devergnette2026sema} are fine-tuned on its training data; earlier pipelines and LLM-assisted annotation \citep{chun2024pipeline,sun2024chinese,gamba2025bootstrapping}. \HAVE

\subsection{Related work}

\paragraph{AMR parsing with language models}
Fine-tuned and instruction-tuned AMR parsers and the limits of prompted LLMs on AMR \citep{lee2023amr,ho2025evaluation,barta2025enhancing,ettinger2023you}.

\paragraph{LLMs as annotators and guideline refinement}
\citet{gilardi2023chatgpt,he2024annollm}; refining guidelines from errors or disagreements on NER \citep{kim2026refining,kim2026diziner}. We differ in the target (a multi-layer graph), the feedback (per layer), and what is revised (instructions, tools, code).

\paragraph{Self-improvement}
\citet{madaan2023selfrefine,shinn2023reflexion,zelikman2024self,lee2026metaharness}; self-correction without external feedback often hurts \citep{huang2024large,kamoi2024when}, which our critic step reproduces (\S\ref{sec:beforeafter}).

% ==================================================================
\section{Method}
\label{sec:method}

\paragraph{Seed annotator}
Sonnet reads only the guideline and writes three steps, \emph{generate}, \emph{critique} and \emph{finalize}, each at most 8,000 characters (17.4K in total for Chinese); every sentence sees the five previous sentences and their annotations, which document-level relations need. No examples, no tools. \HAVE

\paragraph{Feedback per layer}
A judge built on AnCast++ alignments \citep{sun2024anchor,sun2025ancastpp} returns, for each of 25 checks, units with a verdict and the gold and predicted fragments. The builder therefore sees \emph{which layer} each error belongs to. \HAVE

\paragraph{Edits}
In each round the builder sees the current annotator, the training group's gold, every step's output and the incorrect units, and returns operations: edit or rewrite a step's instructions, add a rule step (Python), add or replace a tool. It is told only that PropBank files exist; any lookup it wants, it writes. \HAVE

\paragraph{Admission}
Machine checks only: no copying of test text, at most 10,000 characters per step, tools and rules must pass their own tests; no reviewer model. \HAVE

\paragraph{One round}
Annotate the group, judge it, reflect, edit, re-annotate the group, test on held-out documents. Ten rounds, 50 sentences. Figure~\ref{fig:round} shows Chinese round~3 (temporal anchoring, modal edges on states, \zh{有}\texttt{-01}, coreference of repeated nouns, alignment repair). \HAVE

\begin{figure*}[t]
\centering
\placeholder{0.95\textwidth}{Figure 2: Chinese round 3. Left: incorrect units of the training group (temporal links to DCT, missing modal edges on states, \texttt{have-91} for \zh{有}). Middle: the builder's diagnosis, quoted (``Gold temporal structure is local, not all-DCT''; ``Gold gives states modal edges''; ``Gold uses \zh{有}-01''). Right: the edited lines and the added \texttt{fix\_alignment} rule. Training group 64.4 $\rightarrow$ 82.2 F; test 51.3 $\rightarrow$ 52.6.}
\caption{A round of revision. \HAVE}
\label{fig:round}
\end{figure*}

% ==================================================================
\section{Experimental Setup}
\label{sec:setup}

\paragraph{Data}
Shared-task blind test documents with available gold: 7 Chinese documents (99 sentences) and 4 English (95). Ten training groups of five consecutive sentences: Chinese from ten documents, English from the only three fully annotated documents. The builder never sees test documents. en-0005 and en-0007 are two annotations of the same text; we use them as a human-agreement reference (\S\ref{sec:cannot}). \HAVE

\paragraph{Metrics}
\emph{AnCast++ F} (main): sentence-graph F after \citet{chun2024pipeline}'s cleanup, macro-averaged over documents. \emph{Comprehensive}: adds document-level temporal, modal and coreference relations. \emph{Official \jumach} \citep{stepanek2026first,zeman2026word}: \texttt{compare\_umr.pl} against the shared task's gold, which maps nodes through alignment (the two golds agree at 99.0 on English, 71.7 on Chinese). \emph{Smatch} \citep{cai2013smatch}: sentence graphs against the same gold as AnCast++, duplicate triples removed, unreadable predicted graphs scored as empty, macro-averaged over documents; gold against itself scores 100. Per-layer scores pool the judge's units of a layer. \HAVE

\paragraph{Models and baselines}
Haiku annotates, Sonnet builds. Baselines: the fine-tuned shared-task systems re-scored on our documents; the guideline in the prompt (Haiku, Sonnet), one call per sentence; the guideline read through search tools; one round of guideline-only editing. \HAVE

\paragraph{Reproducibility}
All scores recomputed from released prediction files with the official scorers and our scripts; predictions, every annotator version, reflections, diffs and the edit coding are public. \CHK

% ==================================================================
\section{Results}
\label{sec:results}

\subsection{Reading the guideline is not enough}
\label{sec:diagnosis}

\paragraph{Where the guideline-prompt baseline loses}
Per layer (Table~\ref{tab:layerresults}, ``Prompt''): document-level modality 0.0 / 0.0, document format 11.8 / 10.9, temporal 0.0 / 1.5, coreference 0.0 / 4.4 (Chinese / English), while concepts (65.5 / 82.6) and senses (83.5 / 72.4) are reasonable. The document level is not misunderstood; it is written in the guideline's notation, which the data and the scorers do not use. \HAVE

\paragraph{A Chinese-specific failure}
Chinese concepts score 65.5 because the model translates: of the 742 wrongly labelled concepts whose gold is a Chinese word, 713 (96\%) are English words (\texttt{define-01}); the seed has 170 of 267 and $H_{10}$ 47 of 123. With only the guideline's address and search tools, Chinese falls to 14.9 F and 39.5 Smatch. \HAVE

\paragraph{The seed from the guideline}
The seed fixes part of this on Chinese (F 30.9 $\rightarrow$ 46.4, Smatch 48.0 $\rightarrow$ 59.5) but not on English (F 50.8 $\rightarrow$ 40.9, Smatch 62.8 $\rightarrow$ 54.7): compressing the guideline into checklists keeps its notation and loses detail (\S\ref{sec:beforeafter}). \HAVE

\subsection{Main results}
\label{sec:main}

\begin{table*}[t]
\centering\small
\setlength{\tabcolsep}{4pt}
\begin{tabular}{@{}lcccccccc@{}}
\toprule
& \multicolumn{4}{c}{Chinese} & \multicolumn{4}{c}{English} \\
\cmidrule(lr){2-5}\cmidrule(l){6-9}
System & F & Comp. & \jumach & Smatch & F & Comp. & \jumach & Smatch \\
\midrule
\multicolumn{9}{@{}l}{\emph{Fine-tuned on the shared task's training data (re-scored on our documents)}} \\
Orange & 37.7 & 34.0 & 38.8 & 54.2 & 40.3 & 24.8 & 27.3 & 56.5 \\
Sema & 34.7 & 33.4 & 29.2 & 50.3 & 32.1 & 21.5 & 21.4 & 47.4 \\
\midrule
\multicolumn{9}{@{}l}{\emph{Guideline only}} \\
Guideline in the prompt, Haiku & 30.9 & 24.1 & 12.5 & 48.0 & 50.8 & 29.9 & 21.7 & 62.8 \\
Guideline in the prompt, Sonnet & 50.1 & 38.0 & 24.4 & 64.8 & 54.1 & 31.8 & 28.4 & 65.1 \\
Guideline through tools, Haiku & 14.9 & 12.1 & 8.8 & 39.5 & 46.2 & 27.5 & 19.2 & 59.1 \\
\midrule
\multicolumn{9}{@{}l}{\emph{Using the 50 gold sentences (Haiku)}} \\
Gold sentences as few-shot examples & \tbd & \tbd & \tbd & \tbd & \tbd & \tbd & \tbd & \tbd \\
Guideline edited for 1 round & 32.7 & 30.0 & 36.3 & 43.5$^*$ & 52.5 & 34.3 & 31.8 & 63.7 \\
Guideline edited for 10 rounds & \tbd & \tbd & \tbd & \tbd & \tbd & \tbd & \tbd & \tbd \\
Seed annotator $H_0$ & 46.4 & 34.0 & 24.1 & 59.5 & 40.9 & 22.5 & 25.9 & 54.7 \\
Learned annotator $H_{10}$ & \textbf{52.4} & \textbf{47.2} & \textbf{42.9} & \textbf{65.5} & \textbf{52.2} & \textbf{33.0} & \textbf{38.0} & \textbf{64.6} \\
\bottomrule
\end{tabular}
\caption{Main results on the same test documents. F: AnCast++ sentence-graph F; Comp.: with document-level relations; \jumach: official score against the shared task's gold; Smatch: sentence graphs, macro over documents. $^*$One Chinese document is mostly unreadable in this condition; micro-averaged Smatch is 55.5. TBD rows \RUN.}
\label{tab:main}
\end{table*}

\paragraph{Seed to learned annotator}
Chinese F 46.4 $\rightarrow$ 52.4 (paired document bootstrap 95\% interval 1.0 to 10.1), English 40.9 $\rightarrow$ 52.2 (7.7 to 14.9); Smatch 59.5 $\rightarrow$ 65.5 and 54.7 $\rightarrow$ 64.6; 65 of 98 and 72 of 95 sentences improve by more than a point (Table~\ref{tab:main}). \HAVE

\paragraph{Against the fine-tuned systems}
$H_{10}$ is above Orange on every metric: \jumach 42.9 vs.\ 38.8 and 38.0 vs.\ 27.3, Smatch 65.5 vs.\ 54.2 and 64.6 vs.\ 56.5, from 50 gold sentences instead of the shared task's training data. \HAVE

\paragraph{Against the guideline in the prompt}
On Chinese, the learned Haiku annotator matches Sonnet with the guideline (65.5 vs.\ 64.8 Smatch) and is far above it on the document level (Comp.\ 47.2 vs.\ 38.0, \jumach 42.9 vs.\ 24.4). On English, Haiku with the guideline is already strong on the sentence graph (62.8 Smatch vs.\ 64.6 for $H_{10}$); the gain is on the document level and the official score (Comp.\ 29.9 $\rightarrow$ 33.0, \jumach 21.7 $\rightarrow$ 38.0). \HAVE

\paragraph{Variation between runs}
Re-running $H_0$ and $H_{10}$ with fresh samples changes F by at most 2.1 and Smatch by at most 1.4, below the gains (58.8 / 66.1 Chinese, 56.1 / 63.5 English Smatch). Two more trajectories with other group orders \RUN. \HAVE

\paragraph{Ruling out simpler explanations}
(i) The same 50 gold sentences as few-shot examples: if that is enough, the gain is the data's, not the revision's \RUN. (ii) Editing only the guideline \citep{kim2026refining}: one round fixes notation (\jumach 12.5 $\rightarrow$ 36.3 Chinese) but not structure (F 30.9 $\rightarrow$ 32.7, against 47.2 for the annotator after the same round); ten rounds are needed for a fair comparison \RUN.

\subsection{How many sentences are needed}

\begin{figure}[t]
\centering
\placeholder{0.95\columnwidth}{Figure 3: test F and Smatch of $H_0,\ldots,H_{10}$; training-group F before and after each round; reference lines for the guideline prompt and Orange.}
\caption{Learning curves. \HAVE}
\label{fig:curves}
\end{figure}

\paragraph{The gain comes in the first rounds}
Chinese reaches its full F gain by $H_3$ (52.6 vs.\ 52.4 at $H_{10}$), and Smatch is within 0.3 of $H_{10}$ by $H_3$; English reaches 89\% of its F gain by $H_3$ (51.0) and its best at $H_7$ (53.5). Rounds 4--10 still raise their own training group, by 5--21 F, but not the test documents (Figure~\ref{fig:curves}); \S\ref{sec:analysis} explains why. \HAVE

% ==================================================================
\section{Analysis: What Was Learned and What Was Not}
\label{sec:analysis}

\subsection{The annotator before and after}
\label{sec:beforeafter}

\begin{table*}[t]
\centering\scriptsize
\setlength{\tabcolsep}{4pt}
\begin{tabular}{@{}p{0.07\textwidth}p{0.42\textwidth}p{0.45\textwidth}@{}}
\toprule
Layer & Seed $H_0$ (written from the guideline) & Learned $H_{10}$ (round that introduced it) \\
\midrule
MOD (zh) & \ex{:modstr --- FullAff plain assertion [\ldots]; :modal --- (AUTH :M e) for every top-level event using its modstr} & \ex{GRAPH ATTRIBUTES allowed: :aspect (lowercase), :polarity -, [\ldots] NO :modstr/:modpred/:quot. [\ldots] :modal (author :full-affirmative e) for EVERY event incl.\ state/adjective events} (1, 3) + rule \texttt{root\_modal} (8) \\
CON, SEN (zh) & \ex{Default: Chinese lemma + sense -01 for lexical concepts} & \ex{Ordinary nouns, adjectives-as-modifiers, adverbs, classifiers keep the CHINESE word without sense [\ldots] NEVER translate them. [\ldots] call lookup\_frame(lemma) for every verb [\ldots] NO FRAME -> lemma-00 with :actor and :undergoer} (2) \\
SEN (zh) & \ex{\zh{有} (possession) have-91 (ARG1 possessor, ARG2 possessum)} & \ex{\zh{有} -> \zh{有}-01 (:ARG0 owner/locus, :ARG1 thing had)} (3) \\
ATT (en) & \ex{Common nouns get :ref-number Singular|Plural.} & \ex{Overt plural nouns (men, diseases): NO :refer-number; only pronoun/implicit person nodes get :refer-person + :refer-number.} (2, 5) \\
MOD (en) & \ex{:modpred: the complement of a modal-type predicate (want, [\ldots] think, believe, expect [\ldots])} & \ex{DO NOT write :modstr or :modpred in the graph. [\ldots] (want, need, hope, [\ldots]; NOT think/believe/expect/assume)} (3) \\
SEN (en) & no lookup & \ex{For every event verb/adjective call roleset\_lookup(lemma) to pick the sense}; tool rewritten to read sense names from the frame XML (9) and phrasal rolesets (10) \\
TMP (zh) & \ex{(DCT :depends-on sNt) for relative [time expressions]} & \ex{a past event is (document-creation-time :before e) and two sequential events are (earlier :after later)}; anchoring rule changed in rounds 3, 8 and 9 (\S\ref{sec:cannot}) \\
ASP (zh) & \ex{:aspect --- one of Habitual (usual/repeated), State (stative predicates [\ldots]), Process [\ldots]} & \ex{descriptions of how a work [\ldots] is shot/set/aired/produced (\zh{取景 拍摄 播出 执行}) -> habitual [\ldots] \zh{开始}+V: no \zh{开始} node, V gets :aspect inceptive} (9, 10) \\
\bottomrule
\end{tabular}
\caption{Instruction text of the \emph{generate} step before and after revision (verbatim, abridged). Most learned text replaces the guideline's notation with the data's; the aspect rules are lists of words from training groups.}
\label{tab:beforeafter}
\end{table*}

\paragraph{Size}
Instructions grow from 17.4K to 26.4K characters (Chinese) and 27.9K (English), against 366K for the guideline; Chinese adds 10 rule steps, English 4; each builder writes one PropBank lookup, and the English one is rewritten twice. \HAVE

\paragraph{From the guideline's notation to the data's}
Table~\ref{tab:beforeafter}. The seed writes UMR as the guideline does: modal strength in the graph (\texttt{:modstr FullAff}), \texttt{AUTH}/\texttt{DCT} constants, \texttt{:ref-number Singular} on every noun, \texttt{have-91} for \zh{有}. The learned instructions write it as the data do, and say so explicitly (``NO :modstr'', ``NEVER :refer-number singular on common nouns''). This is the largest single change and it happens in round~1 (English: ``Most of the lost score is format, not understanding'', builder's reflection). \HAVE

\paragraph{The critic had to be restrained}
13 of 84 operations per language restrict the critique and finalize steps (``Critique removed gold-required modals''; ``never flag it''), because the critic ``fixed'' correct annotation. \HAVE

\paragraph{Learned instructions are rules with words, not examples}
$H_{10}$ contains almost no worked examples; it contains rules with short inline lexical examples (\zh{取景 拍摄 播出}; \emph{researchers} $\rightarrow$ \texttt{research-01}). How lexically specific a rule is predicts whether it transfers (\S\ref{sec:cannot}). \HAVE

\subsection{What kinds of edits were made}

\begin{table}[t]
\centering\small
\setlength{\tabcolsep}{4pt}
\resizebox{\columnwidth}{!}{%
\begin{tabular}{@{}llcc@{}}
\toprule
Type & Meaning & zh & en \\
\midrule
CONV & state the data's notation or convention & 38 & 30 \\
CRIT & a decision criterion for the sentence's meaning & 7 & 13 \\
CRITIC & restrain critique / finalize & 13 & 13 \\
TOOLUSE & instruct a step to call a tool & 1 & 3 \\
TOOL & write or replace a tool & 1 & 3 \\
CODE & add a rule step & 10 & 4 \\
SPACE & delete text to stay under the limit & 14 & 18 \\
\midrule
& total applied & 84 & 84 \\
\bottomrule
\end{tabular}}
\caption{Applied operations by type over ten rounds (coding manual and every coded operation in the appendix). \CHK{} (first-pass coding, to be checked by hand)}
\label{tab:edits}
\end{table}

\paragraph{Most edits state conventions}
CONV is the largest type (38 / 30), followed by CRITIC (13 / 13) and space-making deletions (14 / 18); semantic criteria are few (7 / 13), and so are tools (Table~\ref{tab:edits}). The builder spends most of its effort aligning the annotator with how the data are written. \HAVE{} \CHK

\subsection{Which layers moved, and why}
\label{sec:moved}

\begin{table*}[t]
\centering\small
\setlength{\tabcolsep}{3.5pt}
\begin{tabular}{@{}l|rrrrrr|rrrrrrr@{}}
\toprule
& \multicolumn{6}{c|}{Chinese} & \multicolumn{7}{c}{English} \\
Layer & Prompt & $H_0$ & $H_{10}$ & $\Delta$ & Edits & Train & Prompt & $H_0$ & $H_{10}$ & $\Delta$ & Edits & Train & Human \\
\midrule
FMT & 11.8 & 23.0 & 90.5 & +67.5 & 7 & 6.7 & 10.9 & 24.7 & 86.2 & +61.5 & 7 & 1.8 & 89.0 \\
MOD & 0.0 & 0.0 & 49.8 & +49.8 & 12 & 8.3 & 0.0 & 0.0 & 32.0 & +32.0 & 16 & 19.1 & 95.1 \\
ATT & 58.1 & 53.5 & 66.9 & +13.4 & 7 & 21.9 & 15.4 & 11.5 & 53.2 & +41.7 & 10 & 11.3 & 91.6 \\
TMP & 0.0 & 3.0 & 14.7 & +11.7 & 22 & 6.8 & 1.5 & 4.8 & 13.1 & +8.3 & 9 & 3.5 & 90.3 \\
NAM & 47.2 & 44.6 & 56.2 & +11.6 & 3 & 7.9 & 77.6 & 71.4 & 73.3 & +1.9 & 4 & 6.1 & 100.0 \\
CON & 65.5 & 81.8 & 90.4 & +8.6 & 20 & 5.9 & 82.6 & 85.0 & 88.0 & +3.0 & 10 & 4.6 & 99.0 \\
ARG & 56.3 & 60.4 & 66.9 & +6.5 & 15 & 13.0 & 62.6 & 54.9 & 63.2 & +8.3 & 12 & 18.1 & 96.6 \\
SEN & 83.5 & 87.6 & 93.8 & +6.2 & 5 & 4.6 & 72.4 & 66.7 & 80.1 & +13.4 & 20 & 8.7 & 80.3 \\
REE & 25.0 & 30.0 & 35.5 & +5.5 & 8 & 12.3 & 28.5 & 18.9 & 33.0 & +14.1 & 3 & 15.4 & 85.9 \\
COR & 0.0 & 13.7 & 17.3 & +3.6 & 8 & 16.1 & 4.4 & 15.5 & 15.6 & +0.1 & 4 & 2.8 & 91.6 \\
ASP & 66.6 & 74.5 & 74.8 & +0.3 & 9 & 12.8 & 60.9 & 64.3 & 59.6 & $-$4.7 & 10 & 8.2 & 83.4 \\
\bottomrule
\end{tabular}
\caption{Per-layer F on the test documents: guideline-prompt baseline, seed, learned annotator and change. Edits: applied operations targeting the layer (deletions excluded). Train: mean gain of the layer on the training group within a round. Human: en-0007 scored against en-0005, two annotations of the same English text.}
\label{tab:layerresults}
\end{table*}

\paragraph{Conventions transfer almost fully}
The three layers that move most are notation and format: document format +67.5 / +61.5, document modality +49.8 / +32.0, other attributes +13.4 / +41.7 (Table~\ref{tab:layerresults}). Their edits are mostly CONV and CODE, and they move in the round that introduces them and then stay: Chinese format 23.0 $\rightarrow$ 89.7 and modality 0.0 $\rightarrow$ 54.7 in round~1, English attributes 11.5 $\rightarrow$ 40.7 $\rightarrow$ 53.5 in rounds~1--2. A convention seen once applies to every later document. \HAVE

\paragraph{Lookups transfer in part}
Senses improve when the builder gives the model a lookup: Chinese 87.5 $\rightarrow$ 91.9 in round~2 (\texttt{lookup\_frame}) and 94.7 in round~3; English 66.7 $\rightarrow$ 69.9 $\rightarrow$ 78.5 $\rightarrow$ 82.0 over rounds~1--3 (\texttt{roleset\_lookup} in round~1, then rules for phasal verbs and ``deverbal words always take the roleset sense''). With $H_{10}$, 60\% (Chinese) and 39\% (English) of calls use the tool, and removing the tool lowers senses from 94 to 78 (Chinese). English senses reach 80.1, as high as the two human annotations agree (80.3). \HAVE

\paragraph{Rule code fixes format, not meaning}
The 14 rule steps normalize labels, repair variables and alignments, add \texttt{(root :modal author)}, and move modal strength to the document level; removing them does not lower F (Chinese 54.3, English 52.6 without rules), because the instructions learned the same conventions. \HAVE

\subsection{Why some layers do not move}
\label{sec:cannot}

\paragraph{Criteria fit the training group, not the test}
Aspect gets 9 / 10 edits and gains 12.8 / 8.2 points per round on its training group, yet test aspect changes by +0.3 / $-$4.7; coreference gains 16.1 per round in training and +3.6 on test (Chinese). These edits are decision criteria tied to the words of the group (\zh{取景 拍摄 播出 执行} $\rightarrow$ habitual; \zh{开始}+V $\rightarrow$ inceptive; \emph{'d/would}+V $\rightarrow$ habitual), so they cover the cases the builder saw and few others. Contrast format: 6.7 per round in training, +67.5 on test. \HAVE

\paragraph{Temporal anchoring: rules that contradict each other}
Chinese temporal received the most edits (22) and gained only 11.7. Round~3 links only main events to the document creation time (``Gold temporal structure is local''); round~8, trained on a narrative, chains each main event to the previous sentence (test 16.6 $\rightarrow$ 8.6); round~9 reverses it (``the model reinforced its own chaining'', test $\rightarrow$ 14.6). Which events to anchor is a convention that differs between documents, and five sentences of one document show only one of them. \HAVE

\paragraph{Side effects of consolidation}
English aspect falls in rounds~1--2 (64.3 $\rightarrow$ 61.0 $\rightarrow$ 57.7), before any edit targets aspect (first in round~3): rewriting all three steps into the data's notation dropped part of the seed's aspect guidance. Edits are admitted without a held-out check, so such losses stay. \HAVE

\paragraph{Not noisy gold}
On the English text annotated twice, the two annotations agree at 83.4 on aspect, 90.3 on temporal relations and 91.6 on coreference, where $H_{10}$ reaches 59.6, 13.1 and 15.6. The largest gaps to human agreement are exactly the discourse layers whose edits did not transfer; senses, where a lookup exists, reach human agreement. One pair of annotations from one team; they may not be independent. \HAVE{} \CHK

\paragraph{What the builder said about the gold}
Builder suggestions that point to inconsistent gold (e.g.\ \texttt{:refer-number} on plural nouns in some English documents and not in others), counted for the UMR trajectories only \CHK; compared with a human annotator's error notes \ANA.

\subsection{What is learned is about UMR}

\paragraph{Other annotating models}
$H_{10}$, built only from Haiku's errors, run unchanged by Sonnet: Chinese 48.1 $\rightarrow$ 58.9 F and 62.7 $\rightarrow$ 69.8 Smatch over Sonnet's own seed; by GPT: 49.1 $\rightarrow$ 56.4 F and 62.6 $\rightarrow$ 68.3 Smatch. The edits record the scheme's conventions, not one model's weaknesses. Full tables in the appendix. \HAVE

\subsection{Ablations}

\begin{table}[t]
\centering\small
\setlength{\tabcolsep}{3.5pt}
\begin{tabular}{@{}lcccc@{}}
\toprule
& \multicolumn{2}{c}{Chinese} & \multicolumn{2}{c}{English} \\
\cmidrule(lr){2-3}\cmidrule(l){4-5}
Variant & F & Smatch & F & Smatch \\
\midrule
$H_{10}$ & 52.4 & 65.5 & 52.2 & 64.6 \\
seed instructions & 47.4 & 59.9 & 40.0 & 55.0 \\
without tools & 53.0 & 65.4 & 50.6 & 62.5 \\
without rule steps & 54.3 & 66.8 & 52.6 & 63.3 \\
without CONV edits & \tbd & \tbd & \tbd & \tbd \\
without CRIT edits & \tbd & \tbd & \tbd & \tbd \\
\bottomrule
\end{tabular}
\caption{Ablations of $H_{10}$ (tools and rules kept when instructions are swapped). Last two rows \RUN.}
\label{tab:ablation}
\end{table}

\paragraph{The instructions carry the gain}
With the seed's instructions, F falls by 5.0 / 12.3 and Smatch by 5.6 / 9.6; without tools or rule steps it changes by at most 2.1 Smatch (Table~\ref{tab:ablation}). Rolling back one edit type (CONV, CRIT) would divide the gain among types \RUN. \HAVE

% ==================================================================
\section{Conclusion}
\label{sec:conclusion}

\paragraph{Answers}
An LLM with the UMR guideline fails mainly on the project's notation and document level (\S\ref{sec:diagnosis}); fifty gold sentences and an annotator that revises itself are enough to beat the fine-tuned shared-task systems (\S\ref{sec:main}); the gain comes from conventions, which transfer, and lookups, which partly transfer, while criteria for aspect, temporal anchoring and coreference fit the sentences they were learned from (\S\ref{sec:analysis}).

\paragraph{For UMR annotation projects}
The gap between guideline and data is visible from the first annotated sentences and can be closed automatically; the discourse layers need either more data, a resource the model can consult, or people who decide the convention. Next steps: low-resource UMR languages without a PropBank, and a human-in-the-loop annotation tool \citep{anon2026orchestration}.

% ==================================================================
\section*{Limitations}

\paragraph{Scale and variance}
194 test sentences; one trajectory per language (two more in progress); one pair of human annotations for the agreement reference.

\paragraph{Edit coding}
One coder so far; a second coder and agreement are needed.

\paragraph{Models and data}
Builder and main annotator from one model family; possible pretraining contamination; the Chinese AnCast++ gold and the official gold differ.

\paragraph{Design choices}
Prompts, operations, judge and admission rules were designed by the authors.

\section*{Ethics Statement}
Automating annotation can reduce paid annotation work; UMR data are used under their licenses; generative AI was used in experiments and editing.

\bibliography{references}

\appendix

\section{Data and Splits}
Test documents, training windows, the two Chinese golds. \HAVE

\section{The Judge and the Layers}
The 25 checks and how they pool into the eleven layers of Table~\ref{tab:layers}. \HAVE

\section{Prompts}
Output format, seed and reflection prompts. \HAVE

\section{Round-by-Round Records}
Training and test F and Smatch, operations, sizes, per-layer test scores per round. \HAVE

\section{Edit Coding}
The seven types with decision rules, and all 168 coded operations (\texttt{data/annoforge/edit\_codes.tsv}). \HAVE{} \CHK

\section{Other Annotating Models}
Generator/critic pairs and transfer tables, with Smatch. \HAVE

\section{Trajectory Excerpts}
Reflections, the lookup tools and their revisions, rule steps. \HAVE

\section{Cost}
Calls and price per sentence. \HAVE

\end{document}
